% !TeX root = ../navier-stokes-manuscript.tex
% ============================================================
% 2.1 Vortex Stretching
% ============================================================

The momentum equation has four pieces:
\[
  \underbrace{\partial_t u}_{\text{local acceleration}}
  +
  \underbrace{(u\cdot\nabla)u}_{\text{nonlinear transport}}
  =
  \underbrace{-\nabla p}_{\text{pressure}}
  +
  \underbrace{\nu\Delta u}_{\text{viscous diffusion}}.
\]
Nonlinear transport carries the velocity through its own spatial variation, while viscosity diffuses that variation through the surrounding fluid.
The regularity problem turns on whether the smoothing keeps pace wherever the transport sharpens the flow.
The hope is simple: viscosity smooths the fluid faster than nonlinear transport can sharpen it.
If that remains true at every scale, the velocity stays smooth.

\subsection{Vortex Stretching}\label{sub:ThePerfectlyStretchingVortex}

Picture one narrow tube of rotating fluid inside the flow.
Mark off a short piece of that vortex tube and let the same fluid particles form it as time passes.
This is the material segment we will follow: it has an axial length and a cross-section, and both deform as the surrounding velocity field pulls its ends apart.

Write the symmetric part of the velocity gradient, which measures this local deformation, as
\[
  S:=\frac12\bigl(\nabla u+(\nabla u)^{\mathsf T}\bigr),
\]
and let \(e\) be the axial direction and \(L(t)\) the material length.
When the strain has a positive direction aligned with the segment,
\[
  Se=\sigma(t)e,
  \qquad
  \sigma(t)>0,
  \qquad
  \frac{\dot L(t)}{L(t)}
  =
  e\cdot Se
  =
  \sigma(t).
\]
The segment lengthens at the instantaneous logarithmic rate \(\sigma(t)\).

It cannot lengthen while keeping the same cross-section.
Incompressibility preserves the volume of a material region, so the extension along the axis is accompanied by contraction across it.
Let \(\mathcal V\) be the fixed material volume, set \(A(t):=\mathcal V/L(t)\), and write \(r_{\mathrm{eff}}(t):=\sqrt{A(t)/\pi}\).
Then
\[
  \nabla\cdot u=0,
  \qquad
  \frac{d}{dt}(A L)=0,
  \qquad
  \frac{\dot A}{A}=-\sigma(t),
  \qquad
  \frac{\dot r_{\mathrm{eff}}}{r_{\mathrm{eff}}}
  =
  -\frac{\sigma(t)}{2}.
\]
The same positive strain that increases the length decreases the area, and the effective radius falls at half the logarithmic rate.
One motion has made the segment longer and thinner.

The turning carried by the fluid is recorded by the vorticity \(\omega:=\nabla\times u\).
Taking the curl of the momentum equation follows that rotation along the material motion and gives
\[
  \frac{D\omega}{Dt}
  =
  S\omega+\nu\Delta\omega,
  \qquad
  \frac{D}{Dt}
  :=
  \partial_t+u\cdot\nabla,
\]
with the strain term changing the rotation as the segment deforms and the Laplacian recording viscous diffusion of that rotation.
When the vorticity is aligned with the same axial direction, \(\omega=|\omega|e\), the stretching contribution has a definite sign:
\[
  S\omega
  =
  \sigma(t)\omega,
  \qquad
  \frac12\frac{D|\omega|^2}{Dt}
  =
  \sigma(t)|\omega|^2
  +
  \nu\,\omega\cdot\Delta\omega.
\]
On any interval where the right-hand side is positive, the rotation strengthens.
The segment is now lengthening, narrowing, and increasing the rotational part of its velocity gradient at the same time.

The energy carried by the segment depends on its speed and its material volume.
Its volume is already fixed by incompressibility.
If its speed remains bounded, \(|u|\leq U_0\) on \(\Omega_{\mathrm{seg}}(t)\), then
\[
  E_{\mathrm{seg}}(t)
  :=
  \frac12\int_{\Omega_{\mathrm{seg}}(t)}|u(x,t)|^2\,dx
  \leq
  \frac12U_0^2\mathcal V
  <
  \infty,
  \qquad
  \left|\frac12\bigl(\nabla u-(\nabla u)^{\mathsf T}\bigr)\right|_{\mathrm F}
  =
  \frac{|\omega|}{\sqrt2}
  \leq
  |\nabla u|_{\mathrm F}.
\]
Bounded speed and fixed material volume keep the segment energy finite, while \(|\nabla u|_{\mathrm F}\) grows wherever \(|\omega|\) grows.
On such an interval, the cross-section is shrinking while the rotating velocity gradient concentrates inside the same material tube.
To follow that narrowing into the equation, we compare the original field with a rescaled solution at a finer width.
The rescaled field is not a later state of this material tube.

\divider{\NavierStokes{} Scaling}

Fix a time \(t_0\) at which the original segment has effective width \(\ell:=r_{\mathrm{eff}}(t_0)\).
For \(\lambda>1\), compare it at the same fixed viscosity with the rescaled fields
\[
  u_\lambda(x,t)
  :=
  \lambda u(\lambda x,\lambda^2t),
  \qquad
  p_\lambda(x,t)
  :=
  \lambda^2p(\lambda x,\lambda^2t).
\]
At time \(\lambda^{-2}t_0\), the corresponding vortex in \(u_\lambda\) has width \(\lambda^{-1}\ell\).
The comparison lets the equation reveal what changes when the width is reduced.
Differentiating the rescaled fields gives
\[
\begin{aligned}
  \partial_tu_\lambda(x,t)
  &=
  \lambda^3(\partial_tu)(\lambda x,\lambda^2t),
  \\
  (u_\lambda\cdot\nabla)u_\lambda(x,t)
  &=
  \lambda^3\bigl((u\cdot\nabla)u\bigr)(\lambda x,\lambda^2t),
  \\
  \nabla p_\lambda(x,t)
  &=
  \lambda^3(\nabla p)(\lambda x,\lambda^2t),
  \\
  \nu\Delta u_\lambda(x,t)
  &=
  \lambda^3\nu(\Delta u)(\lambda x,\lambda^2t).
\end{aligned}
\]
Every term acquires the same factor \(\lambda^3\).
Viscous diffusion is stronger in the finer comparison, but nonlinear transport is stronger by the same amount, with pressure and local acceleration retaining their place in the equation.
Reducing the width does not tilt the balance toward viscosity.

\divider{Critical Height}

The four terms keep their relative strength under this concentration, but the familiar sizes of the velocity do not.
Let \(\Lambda:=(-\Delta)^{1/2}\).
A change of variables gives
\[
  \norm{u_\lambda(t)}_{\dot H^s}^2
  =
  \lambda^{2s-1}
  \norm{u(\lambda^2t)}_{\dot H^s}^2.
\]
There is one height at which the power of \(\lambda\) vanishes.
At \(s=\tfrac12\), set
\[
  H(t)
  :=
  \frac12\norm{u(t)}_{\dot H^{1/2}}^2
  =
  \frac12\norm{\Lambda^{1/2}u(t)}_2^2.
\]
Writing \(H_\lambda\) for the same quantity evaluated on \(u_\lambda\), the energy, critical height, and first-derivative energy scale side by side:
\[
  \norm{u_\lambda(t)}_2^2
  =
  \lambda^{-1}\norm{u(\lambda^2t)}_2^2,
  \qquad
  H_\lambda(t)=H(\lambda^2t),
  \qquad
  \norm{\nabla u_\lambda(t)}_2^2
  =
  \lambda\norm{\nabla u(\lambda^2t)}_2^2.
\]
Kinetic energy falls as the first-derivative energy rises, while the critical height does not move.
It is the scale at which the concentration remains fully visible without the norm itself being enlarged or diminished by rescaling.

\divider{Nonlinear Production versus Viscous Dissipation}

The invariant height locates the balance, and its evolution along the original solution shows which side is increasing it.
Apply \(\Lambda^{1/2}\) to the momentum equation and pair with \(\Lambda^{1/2}u\).
Write the signed nonlinear production as
\[
  \mathcal P_{1/2}(t)
  :=
  -\operatorname{Re}
  \pair
    {\Lambda^{1/2}u(t)}
    {\Lambda^{1/2}\bigl((u(t)\cdot\nabla)u(t)\bigr)}_{L^2}.
\]
Incompressibility makes the pressure pairing vanish, and the Laplacian contributes \(-\nu\norm{\Lambda^{3/2}u}_2^2\).
The critical height evolves by
\[
  H'(t)
  +
  \nu\norm{\Lambda^{3/2}u(t)}_2^2
  =
  \mathcal P_{1/2}(t).
\]
The height rises when nonlinear production exceeds viscous dissipation.
Under the finer-scale comparison, these two contributions become
\[
  \mathcal P_{1/2}[u_\lambda](t)
  =
  \lambda^2\mathcal P_{1/2}[u](\lambda^2t),
  \qquad
  \nu\norm{\Lambda^{3/2}u_\lambda(t)}_2^2
  =
  \lambda^2\nu\norm{\Lambda^{3/2}u(\lambda^2t)}_2^2.
\]
They remain matched by the common factor \(\lambda^2\).
That is also the inverse of the time contraction \(t\mapsto\lambda^{-2}t\): a finer spatial scale comes with a faster response scale.

\divider{The Heat Clock}

The square relation between width and response time can be read directly from viscosity.
Set the other terms aside for this comparison and follow the linear heat equation \(v_t=\nu\Delta v\).
A Fourier mode evolves according to
\[
  \widehat v(\xi,t+\delta t)
  =
  e^{-\nu|\xi|^2\delta t}\widehat v(\xi,t).
\]
A structure of width \(r\) is concentrated near frequencies \(|\xi|\simeq r^{-1}\).
Viscosity changes it by order one when the exponent \(\nu|\xi|^2\delta t\) is of order one, which gives the comparison time
\[
  \tau_\nu(r)
  :=
  \frac{r^2}{\nu}.
\]
This is the heat clock attached to width \(r\).
At the finer comparison width,
\[
  \tau_\nu(\lambda^{-1}r)
  =
  \lambda^{-2}\tau_\nu(r),
\]
so each spatial contraction shortens the viscous comparison time by the same square factor carried by the full \NavierStokes{} scaling.
Repeated concentration now shortens the associated comparison time at every step.

\divider{The Zeno Cascade}

Repeat the contraction geometrically.
For the widths and their heat clocks, set
\[
  r_n:=2^{-n}r_0,
  \qquad
  \tau_n:=\frac{r_n^2}{\nu},
\]
so each width is half the one before it, while each clock is one quarter as long.
The complete sequence of comparison times is summable:
\[
  \tau_n
  =
  \frac{r_0^2}{\nu}4^{-n},
  \qquad
  \sum_{n=0}^{\infty}\tau_n
  =
  \frac{4r_0^2}{3\nu}
  <
  \infty.
\]
This does not construct a cascade or establish a singularity.
It shows that the heat clocks alone cannot rule out infinitely many scale transitions before a finite time.

The rescaled fields were comparisons.
For those clocks to describe an actual candidate singularity, one continuing \NavierStokes{} solution would have to pass through the widths
\[
  r_0>r_1>r_2>\cdots\longrightarrow0
\]
at times
\[
  t_0<t_1<t_2<\cdots\uparrow T_*<\infty,
  \qquad
  t_{n+1}-t_n\asymp\tau_n,
\]
with comparison constants independent of \(n\).
Its remaining time at the \(n\)-th width would then satisfy
\[
  T_*-t_n
  \asymp
  \sum_{k=n}^{\infty}\tau_k
  =
  \frac{4r_n^2}{3\nu}.
\]
Each finer structure would have to form inside an interval shrinking at the same rate as its heat clock.
The spatial sequence has become a demand on how rapidly the same velocity field can change.
Following that demand requires its velocity, acceleration, jerk, and every later temporal response.
